%!TEX program = lualatex
\documentclass[ngerman,headheight=70pt]{scrartcl}
\usepackage[ngerman]{babel}
\usepackage{blindtext}
\usepackage[utf8]{luainputenc}
\usepackage{scrlayer-scrpage}
\usepackage[T1]{fontenc}
\usepackage{fontspec}
\usepackage{multicol}
\usepackage{enumitem}
\usepackage{epigraph}
\usepackage[german=quotes]{csquotes}
\usepackage[bottom=2cm,footskip=8mm]{geometry}
\setmainfont{Source Sans Pro}
\MakeOuterQuote{"}

\setlength{\parindent}{0pt}
\setlength{\parskip}{6pt}

\pagestyle{plain.scrheadings}
\thispagestyle{scrheadings}
\lohead{\vspace*{-1.5cm}\includegraphics[scale=0.4]{../images/stupa_logo_rot.pdf}}
\rohead{\\Präsidium des\\Studierendenparlaments}

\lofoot[Universität Hamburg · Präsidium des Studierendenparlaments\\\hrule Von-Melle-Park 5 · D-20146 Hamburg]{Universität Hamburg · Präsidium des Studierendenparlaments\\\hrule Von-Melle-Park 5 · D-20146 Hamburg}
\cofoot[]{}
\cohead[\thepage]{}

\renewcommand{\thesubsection}{\arabic{section}.\arabic{subsection}}
\renewcommand{\thesubsubsection}{\arabic{subsubsection}.}
\setcounter{subsection}{-1}

\begin{document}
    UHH · StuPa-Präsidium · Von-Melle-Park 5 · D-20146 Hamburg

    \section{Inhaltliches Konzept für das Wintersemester}

    \subsection{StuPa-Sitzungen}

    Am 27. Oktober werden wir als Präsidium die Themen "Wahlbeteiligung bei der
    StuPa-Wahl" und "Haushaltsberatungen in der Bürgerschaft" als TOPs vorschlagen.

    Der Hintergrund ist die geringe Wahlbeteiligung von rund 13\% bei der letzten
    Wahl und wir möchten diese Gelegenheit nutzen, um Ideen zu sammeln, was zur
    Erhöhung der Wahlbeteiligung gemacht werden kann. Jede*r ist herzlich eingeladen
    sich daran zu beteiligen.

    Das andere Thema ergibt sich aus dem Fakt, dass zu dem Zeitpunkt der nächste
    Haushalt in den Fachausschüssen beraten wird und Anfang November der
    Haushaltsausschuss die Stellungnahmen zusammentragen wird. Die Oktobersitzung
    ist daher die einzige realistische Möglichkeit, um noch einen Antrag zu
    verabschieden, der eine Ausfinanzierung der Uni oder aber mindestens den
    Ausgleich der Lohnsteigerungen fordert und diesen dann auch noch den
    Abgeordneten zukommen zu lassen.

    Bis zum 27. Oktober ist es unsere Absicht die Faktenlage zu sichten und
    herauszubekommen, inwieweit sich das Budget der Uni im Vergleich zum letzten
    Haushalt verändert hat, um auf der Grundlage spezifische Veränderungen oder
    Anpassungen fordern zu können.

    Angesichts der Tatsache, dass eine Unterfinanzierung die Universität als Ganzes
    bedroht und bereits einige Studiengänge gefährdet sind oder sogar bereits
    nicht mehr angeboten werden, wäre es daher wünschenswert, wenn die Verfasste
    Studierendenschaft gemeinsam für eine Verbesserung der Situation eintritt.
    Wenn auch die Oppositionsfraktionen einem solchen Antrag zustimmen, wäre
    die Wirkmächtigkeit natürlich weitaus höher.

    Am 15. Dezember werden wir als Präsidium vorschlagen z.B. im Rahmen einer
    Aktuellen Stunde über die Haushaltsberatungen zu sprechen, denn am 13., 14.
    und 15. Dezember sind die Haushaltsberatungen in der Bürgerschaft an deren
    Ende der Haushalt verabschiedet wird.

    \subsection{Wahlmotto}

    Nach reiflicher Überlegung sind wir zum Wahlmotto für die diesjährige Wahl
    des Studierendenparlaments gekommen:

    "Fremdbestimmt oder selbstgestaltet? Universität gemeinsam entwickeln!"

    \subsection{Wahlveranstaltungen}

    Noch vor Beginn der heißen Wahlkampfphase werden wir eine inhaltliche Veranstaltung
    zum Thema Studienreform veranstalten. Insbesondere soll dort über die Frage
    "Wie soll Studium aussehen?" diskutiert werden. Diese Veranstaltung wird am
    10. November stattfinden.

    Vor der StuPa-Sitzung am 15. Dezember wird es die erste Listenvorstellung
    geben, wobei das Konzept dieses Jahr gänzlich anders ist. Die erste
    Vorstellung wird sich um die Weltanschauung der Listen drehen.
    Dabei wird es um die Fragen "Wie sieht Hamburg im Jahr 2030 aus?",
    "Wie soll Gesellschaft aussehen?" und "Wie soll Hochschule aussehen?"
    gehen. Das Ziel der ersten Vorstellung ist somit über die grundsätzlichen
    Ansichten einer Hochschulgruppe zu reden statt über konkrete Detailfragen.

    Vor der StuPa-Sitzung am 12. Januar gibt es die zweite Listenvorstellung,
    bei der es um die Frage "Was wollt ihr in einem Jahr in der Verfassten
    Studierendenschaft machen?" gehen wird.

    Diese beiden Vorstellungen geben den potentiellen Wähler*innen, wie den
    Hochschulgruppen die Möglichkeit einmal aus dem (Gremien-)alltag aus- und in
    eine viel zu selten geführte Diskussion über die Visionen und Motive einzutreten.

    Wie in den vergangenen Jahren soll es die Möglichkeit geben die Liste im
    Vorfeld der Veranstaltungen in Form von Infoständen vorzustellen.

\end{document}
